







\documentclass[11pt]{article}
\usepackage{latexsym}
%\usepackage{showlabels}
%\setlength{\textheight}{7.5in}
%\setlength{\textwidth}{5.2in}
%\flushbottom
%\documentclass[10pt]{article}
\usepackage{amsmath,amsthm,amsfonts,amssymb,graphicx}
\usepackage{comment}
%\usepackage{showkeys}
%\usepackage{epsf}
%\raggedbottom
%\textwidth 6in
%\oddsidemargin .25in
%\evensidemargin .25in
%\textheight 8.5in
%\topmargin 0in





\newtheorem{thm}{Theorem}[section]
\newtheorem{q}[thm]{Q}
\newtheorem{lemma}[thm]{Lemma}
\newtheorem{cor}[thm]{Corollary}
\newtheorem{prop}[thm]{Proposition}
\newtheorem*{main}{Main Theorem}{}

\theoremstyle{remark}
\newtheorem{example}[thm]{Example}
\newtheorem{exercise}[thm]{Exercise}
\newtheorem{definition}[thm]{Definition}
\newtheorem{remark}[thm]{Remark}
\newtheorem{examples}[thm]{Examples}

\def\b{\begin{q}}
\def\e{\end{q}}


\bibliographystyle{amsalpha}   %plain


\def\R{{\mathbb R}}
\def\Z{{\mathbb Z}}
\def\H{{\mathbb H}} 
\def\Q{{\mathbb Q}} 
\def\G{{\Gamma}}
\newcommand{\<}{\langle}
\renewcommand{\>}{\rangle}

\begin{document}
\centerline{\Huge\bf Questions in Geometric Group Theory}
\vskip 1cm
\centerline{\Large Collected by Mladen Bestvina}
\centerline{\tt bestvina@math.utah.edu}
\centerline{\tt http://www.math.utah.edu/$\sim$bestvina}
\centerline{Major revision August 22, 2000}
\vskip 1cm
Notes: 


I regret that I am unable to offer 1M for solutions.

Disclaimer: Names in parentheses reflect the person I heard the question from.
The actual source may be different. Any corrections, new questions, additional
references, clarifications, solutions and info welcome.

For questions with a decidedly more combinatorial flavor, see \newline
{\tt
  http://www.grouptheory.org}

\section{Hyperbolic Groups}

\subsection{Characterization and Structure}



\begin{q}
Suppose $G$ admits a finite $K(G,1)$. If $G$ does not contain any
Baumslag-Solitar subgroups $BS(m,n)$, is $G$ necessarily hyperbolic? If $G$
embeds in a hyperbolic group, is it hyperbolic? 
\end{q}

The answer is no if the
``finite $K(G,1)$'' assumption is replaced by ``finitely presented'' by
[Brady].

Remarks: $BS(m,n)=<x,y|x^{-1}y^mx=y^n>$ for $m,n\neq 0$.
One considers $\Z\times\Z$ as a Baumslag-Solitar group $BS(1,1)$. 

The answer is not known even when $Y=K(G,1)$ is a non-positively curved
2-complex. In that case, the question becomes: If the universal cover $\tilde
Y$ contains a flat, does $G$ contain $\Z\times\Z$?

It is possible that there are counterexamples that contain certain infinite
quotients of Baumslag-Solitar groups. A modification of the above question
would be: 

\begin{q}
Suppose $G$ admits a finite $K(G,1)$, does not contain $\Z\times\Z$, and
whenever $x\in G$ is an infinite order element such that $x^m$ and $x^n$ are
conjugate, then $|m|=|n|$. Is $G$ hyperbolic?
\end{q}



All known constructions of hyperbolic groups of large (rational, or integral
for torsion-free groups) cohomological dimension use arithmetic groups as
building blocks. See [Charney-Davis, Strict Hyperbolization], [Gromov,
Asymptotic Invariants, Ch.7]. Gromov loosely conjectures that this is always
the case. The following is an attempt to make the question precise.

\begin{q}
For every $K>0$ there is $N>0$ such that every
    word-hyperbolic group $G$ of rational cohomological dimension $\geq N$
    contains an arithmetic lattice of dimension $\geq K$.
\end{q}

Note: An earlier formulation of this question required $N=K$ for $K$ large. A
counterexample to the stronger formulation was found by [Mosher-Sageev], see
{\tt http://andromeda.rutgers.edu$/\sim$mosher/HighDHyp.ps}. See also [Gromov,
Asymptotic Invariants, 7.VI]. One might want to allow certain quotients of
arithmetic groups as well.

The following is another instance of Gromov's loose conjecture.

\begin{q}
Conjecture (Davis): All hyperbolic
Coxeter groups have uniformly bounded rational (or virtual) cohomological
dimension. 
\end{q}
[Moussong] has a criterion for recognizing word-hyperbolicity
from the Coxeter diagram.

Update (May 2002): T. Januszkiewicz and J. Swiatkowski have
constructed hyperbolic Coxeter groups of arbitrarily large
dimension. This answers both questions above. The preprint is available
at

{\tt http://www.math.uni.wroc.pl/$\sim$tjan/Papers/submitted.html}

\begin{q}
(Davis) If $G$ is word-hyperbolic, does the Rips complex $P_d(G)$
have an equivariant negatively curved metric for $d$ sufficiently large?
\end{q}

A potential counterexample is the mapping torus of a hyperbolic automorphism
of a free group. 

\subsection{Subgroups of Hyperbolic Groups}

\begin{q}
(Gromov) Does every 1-ended word-hyperbolic group contain a
closed hyperbolic surface subgroup?
\end{q} 

\begin{q}
(Gromov) For a given $n$ is there an example of a hyperbolic group of
dimension $n$ in which every infinite index subgroup is free? Or in which
there are no (quasi-convex) subgroups with codimension $\leq k$ for a given
$k\leq n-2$.
\end{q}



\begin{q}\label{conjugates}
  (Swarup) Suppose $H$ is a finitely presented subgroup of a word-hyperbolic
  group $G$ which has finite index in its normalizer. Assume that there is
  $n>0$ such that the intersection of $n$ distinct conjugates of $H$ is always
  finite. Is $H$ quasi-convex in $G$?
\end{q}

The converse is a theorem of
[Gitik-Mitra-Rips-Sageev]. A special case worth considering is when
$G$ splits over $H$ when Gersten's converse of the combination theorem
might be helpful.

Remark(Gitik): The problem is open even when $H$ is malnormal in $G$.

\begin{q}
(Mitra) Let $X_G$ be a finite 2-complex with fundamental group $G$. Let
$X_H$ be a cover corresponding to the f.p. subgroup $H$. Let $I(x)$
denote the injectivity radius of $X_H$ at $x$.
Does $I(x) \rightarrow \infty$ as $x\rightarrow\infty$
imply that $H$ is quasi-isometrically embedded in $G$ ?
\end{q}
 
A positive answer to the above question for $G$ hyperbolic would imply
a positive answer to Q \ref{conjugates}.

\b
(Canary) Let $G$ be word-hyperbolic and $H$ a f.p. subgroup of $G$.
Suppose that for all $g\in G$ there is $n>0$ such that $g^n\in H$. Does
it follow that $H$ has finite index in $G$? \e

Yes if $H$ is quasi-convex,
since then $\Lambda(H)=\Lambda(G)$.

\b
(Whyte)
Let $\Gamma$ be a 1-ended hyperbolic group. Can every infinite index
subgroup be free?
\e

\b
(Whyte)
Let $\Gamma$ be a 1-ended hyperbolic group. Can a finite index
subgroup of $\Gamma$ be isomorphic to a subgroup of $\Gamma$ of
infinite index?
\e

\subsection{Relative Questions}

\begin{q}
(Swarup) Prove the combination theorem for relatively hyperbolic
groups.
\end{q}



\b
(Swarup) There is a theorem of Bowditch saying that if $G$
acts on a compact metrizable space $X$ so that the action is 
properly discontinuous and cocompact on triples, then $G$ is a hyperbolic
group and $X$ is its boundary. State and prove the analog in which one
allows parabolics.
\e









\subsection{Residual Finiteness}

\b
Is every word-hyperbolic group residually finite?
\e

\noindent Notes: 

a) D. Wise has constructed a finite 2-dimensional locally CAT(0) complex
whose fundamental group is not RF.

b) Z. Sela showed that torsion-free word-hyperbolic groups are Hopfian.

c) M. Kapovich pointed out that there are non-linear word-hyperbolic
groups.  Start with a lattice $G$ in quaternionic hyperbolic space and
adjoin a ``random'' relation to get the desired group. Super-rigidity
[Corlette, Gromov-Schoen]
says that every linear representation of $G$ is either faithful or has
finite image.

\b
(Dani Wise)
{\rm Let $G^n$ denote the Cartesian product of $n$ copies of the 
group $G$, and let $rank(G^n)$ be the smallest number of generators of $G^n$.} 

\noindent
{\bf Conjecture.} If $G$ is word-hyperbolic then
$\lim_{n\to\infty}rank(G^n)=\infty$.\e

The Conjecture is true if $G$ is finite and nontrivial, as can be seen by a
simple pigeon-hole argument. It is also true if $G$ has a quotient for which
the conjecture holds. In particular, it is true for groups that have a
proper finite index subgroup. 

The Conjecture is false if there is an epimorphism $G\to G\times G$. There is
an example (Wise) of a 2-generator infinitely presented $C'(\frac{1}{6})$
small cancellation group where the conjecture fails.

\b (Dani Wise)
Find ``nice'' (e.g. $CAT(0)$, automatic,...) groups where this
conjecture fails.\e

\subsection{Algorithms}

\b (Epstein) Let $G$ be a word-hyperbolic group and $\partial G$ its boundary.
Is there an algorithm to compute $\mathaccent20{H}^i(\partial G)\cong
H^{i+1}(G,\Z G)$? In particular, is there an algorithm to decide whether
$\mathaccent20{H}^i(\partial G)\cong \mathaccent20{H}^i(S^2)$ for all $i$? 
\e
If
$\partial G$ has the cohomology of $S^2$ then it is homeomorphic to $S^2$
(Bestvina-Mess) and modulo a finite normal subgroup $G$ is conjecturally
commensurable to a hyperbolic 3-manifold group. 

{\it Remark} (Epstein, Sela) There is an algorithmic procedure to determine
the number of ends (i.e. $\mathaccent20{H}^0(\partial G)$) of a hyperbolic
group. First, one finds algorithmically an explicit $\delta$ (of
$\delta$-hyperbolicity). Then one finds an automatic structure, from which it
can be immediately read if the group is finite or 2-ended. One can run an
(enumeration) machine that terminates if the group nontrivially splits over a
finite subgroup, i.e. if it has infinitely many ends. Finally, there is a
machine (Gerasimov) that terminates if the group is 1-ended. This is based on
the property of 1-ended groups (proved by Bowditch) that there is a universal
bound on the length of a shortest path connecting points on a sphere $S(R)$ at
distance $<100+100\delta$ and disjoint from the sphere $S(R-10-10\delta)$.

{\it Remark} (Sela)
For 1-ended torsion-free hyperbolic groups there is an algorithm
based on Sela's work on the isomorphism problem to decide if the group splits
over $\Z$. This algorithm can be ``relativized'' to find the JSJ decomposition
as well. The details of this have not appeared. The case of groups with
torsion is open. 



\subsection{Maps Between Boundaries}

\b
(M. Mitra) Let $G$ be a word-hyperbolic group and $H$ a
word-hyperbolic subgroup. Does the inclusion $H\to G$ extend to a
continuous map between the boundaries $\partial H\to\partial G$? \e

A theorem of Cannon-Thurston says that this is the case when $G$ is the
fundamental group of a hyperbolic 3-manifold that fibers over $S^1$ and $H$ is
the group of the fiber. The map $S^1\to S^2$ is surjective and it can be
explicitely described in terms of the stable and unstable laminations of the
monodromy. When the manifold is closed, the map is finite-to-1.

By a theorem of Mitra, if $G$ is a graph of groups with one vertex group $H$,
all vertex groups hyperbolic, and all edge-to-vertex monomorphisms
quasi-isometric embeddings, then the answer is yes.

\b
(Swarup) Suppose $G$ is a hyperbolic group which is a graph of
hyperbolic groups such that all edge to vertex inclusions are
quasi-isometric embeddings.  [Mitra] shows that for each vertex group
$V$ inclusion $V\hookrightarrow G$ induces a continuous
Cannon-Thurston map $\partial V\to \partial G$.  Describe the
point-preimages. In particular, show that the map is finite-to-one.
\e



\subsection{Miscellaneous}

\b \label{thurston}
(Thurston) Is every closed hyperbolic 3-manifold finitely covered
by one that fibers over the circle?\e

\b
(Jim Anderson) Can the fundamental group of a hyperbolic 3-manifold
that fibers over the circle contain a subgroup that is locally free but
not free? \e
There are such hyperbolic 3-manifolds, so the negative answer
to this question would provide a counterexample to Q \ref{thurston}.

Richard Kent constructs an infinite family of examples. See\newline
{\tt http://www.math.utexas.edu/$\sim$rkent/cuff}

\b (Ian Leary)
Is there a version of the Kan-Thurston theorem using only 
CAT(-1) groups, or word hyperbolic groups?  (The statement 
should be:  for any finite simplicial complex X, there is a
locally CAT(-1) polyhedral complex $Y$ and a map $Y \to X$ 
that is surjective on fundamental groups and induces an 
isomorphism on homology for any local coefficients on $X$.)  
\e

\b (Ian Leary)
Is there any restriction on the homotopy type of the 
quotient $R_d(G)/G$ for $G$ word hyperbolic?  
\e
Here $R_d$ is 
the Rips complex, and $d$ should be taken to be large 
(so that $R_d(G)$ is a model for the universal proper $G$-space).  
     Of course this space is a finite complex, but is there any 
other restriction on its homotopy type?   

Background: Leary-Nucinkis proved that any complex has the homotopy 
type of $\underline EG/G$ for some discrete group $G$, where
$\underline EG$ denotes 
the universal proper $G$-space.  (This is a version of the
Kan-Thurston 
theorem for $\underline EG$ instead of $EG$.)  

Leary can answer the analogues of both questions above with CAT(0) 
in place of CAT(-1) or word hyp., as well as the 2-dimensional 
cases of both using either CAT(-1) groups or small cancellation 
groups. The trouble with getting to higher dimensions is that 
most proofs of Kan-Thurston type results use direct products to 
make higher dimensional groups.  

\section{$CAT(0)$ groups}


\b
(Swarup) Is there a proof of Johannson's theorem that $Out(\pi_1
M)$ is virtually generated by Dehn twists for $M$ a Haken 3-manifold
along the lines of Rips-Sela's theorem that $Out(G)$ is virtually
generated by Dehn twists for hyperbolic $G$? Is this true for CAT(0)
groups? In particular, if $G$ is a CAT(0) group and $Out(G)$ is
infinite, does $G$ admit a Dehn twist of infinite order? 
\e
By a Dehn
twist we mean an automorphism $\phi$ of the following form: either $G$
splits as $A*_CB$, $t\in C$ is central, $\phi(a)=a$ for $a\in A$ and
$\phi(b)=t^{-1}bt$ for $b\in B$; or $G$ is an HNN extension and $\phi$
is defined similarly.



\b
(Gromov) If $G$ admits a finite dimensional $K(G,1)$, does
$G$ act properly discontinuously by isometries on a complete CAT(0)
space?
\e

\b
(Eilenberg-Ganea) Is there a group $G$ of cohomological dimension 2 and
geometric dimension 3?\e

\b
(Whitehead) Is every subcomplex of an aspherical 2-complex aspherical?\e



\b
(Exercise in [Ballmann-Gromov-Schroeder, p.2]) Take a closed
surface $S$ of genus $\geq 2$. Let $V=S\times S$ and let
$\Sigma\subset V$ denote the diagonal. Let $\tilde V$ be a
nontrivially ramified finite cover of $V$ along $\Sigma$. Then $\tilde
V$ has a natural piecewise hyperbolic CAT(0) metric. Show that $\tilde
V$ admits no $C^2$-smooth Riemannian metric with curvature $K\leq 0$.\e

\b
Suppose a group $G$ acts properly discontinuously and cocompactly
by isometries on two CAT(0) spaces $X$ and $Y$. Croke-Kleiner have
examples where the boundaries $\partial X$ and $\partial Y$ are not
equivariantly homeomorphic. Is there a compact metric space $Z$ and
cell-like maps $Z\to\partial X$, $Z\to\partial Y$? \e
A surjective map
between metric compacta is {\it cell-like} if each point preimage is
cell-like. A compact metric space is cell-like if when embedded in the
Hilbert cube $I^\infty$ (or $I^n$ if it is finite-dimensional) it is
contractible in each of its open neighborhoods.

\b
(D. Wise) Let $G$ act properly discontinuously 
and cocompactly on a CAT(0) space
(or let $G$ be automatic).
Consider two elements $a,b$ of $G$.
Does there exist $n>0$ such that 
either the subgroup
$\langle a^n,b^n\rangle$ is free
or
$\langle a^n,b^n\rangle$ is abelian?\e

\b
Do CAT(0) (or (bi)automatic) groups satisfy the Tits
alternative?\e

\b Does every Artin group have a finite $K(G,1)$? \e
Yes for Artin
groups of finite type (meaning that the associated Coxeter group is
finite) by the work of Deligne.   

\b (Bestvina)
Let $L_1,L_2,\cdots,L_n$ be a finite collection of straight lines in the
plane $\R^2$. Does the fundamental group of $$X=\R^2\times\R^2\setminus
\bigcup_{i=1}^n L_i$$ admit a finite $K(\pi,1)$?\e

By [Deligne] the answer is yes for simplicial arrangements. In general $X$ is
not aspherical, but in all examples one understands, $X$ becomes aspherical
after attaching finitely many $\geq 3$-cells. It is not even known whether
$\pi_1(X)$ is torsion-free. A similar question can be asked about
finite hyperplane arrangements in $\R^m$.

\b
(Eric Swenson)
Let $X$ be a $CAT(0)$ metric space and $G$ a finitely generated group that
acts properly discontinuously by isometries on $X$.

1. Can $G$ be an infinite torsion group? The expected answer is yes.

2. If the action is cocompact, can $G$ contain an infinite torsion subgroup?
   The expected answer is no.
\e

\b
(Kim Ruane)
Let $G$ be a Coxeter group, e.g. right-angled,  and assume that $G$ acts
properly discontinuously and by isometries on a $CAT(0)$ space $X$. How is $X$
different from the Coxeter complex? Specifically:

Suppose that $H$ is a special subgroup of $G$. Is there a closed convex
subset of $X$ on which $H$ acts cocompactly?\e


\b (Ruth Charney) Classify Coxeter groups up to isomorphism. \e
Interesting examples of isomorphic Coxeter groups (and Artin groups)
with non-isomorphic diagrams were given by Noel Brady, Jon McCammond,
Bernhard Muehlherr, and Walter Neumann. In
the opposite direction, conditions under which isomorphism of groups
implies isomorphism of diagrams were given by Charney-Davis, David
Radcliffe, Anton Kaul.  \b (Ruth Charney) Classify Artin groups up to
isomorphism. \e

\b
(Ruth Charney)
Are all finite type Artin groups linear? \e
Braid groups are linear by the
   recent work of Bigelow and Krammer. 

Update: A.M. Cohen and D.B. Wales and independently F. Digne show that
the answer is yes. See {\tt
http://front.math.ucdavis.edu/math.GR/0010204} and {\tt
http://www.mathinfo.u-picardie.fr/digne} 

\b (Ruth Charney) Are all
[finite type] Artin groups $CAT(0)$? \e The answer is yes for small
numbers of generators by the work of Krammer, Tom Brady, Jon
McCammond. The question is open even for braid groups.

\b
(Ruth Charney)
Are all Artin groups automatic?\e

\b
(Ross Geoghegan)
Let $M$ be a proper $CAT(0)$ space. We say that $M$ is {\it almost
  geodesically complete} if there is $R\geq 0$ such that for all $a,b\in M$
  there is an infinite geodesic ray starting at $a$ and passing within $R$ of
  $b$. 

If $Isom(M)$ acts cocompactly on $M$, is $M$ almost
geodesically complete? \e

Note: It is a theorem of Ontaneda that if a discrete group acts cocompactly by
isometries on $M$ and if $H^*_c(M)\neq 0$ then $M$ is almost geodesically
complete. 

\b
(Dani Wise)
{\rm A {\em triplane} is a CAT(0) space obtained by gluing together
three Euclidean half-planes along their boundaries.
%
 Say that a CAT(0) space $X$ has {\em isolated flats} if $X$ does not
contain an isometrically embedded triplane.

Let $G$ act on a CAT(0) space $X$. Recall that a subgroup $H$ is
quasiconvex relative to this action provided that for a point
$x\in X$ there is a constant $K$ such that for any two points in
the orbit $Hx$, the geodesic connecting these points lies in a
$K$-neighborhood of $Hx$.}

{\bf Conjecture:} Let $G$ act properly discontinuously and
cocompactly on a CAT(0) space $X$ with isolated flats. Let $H$ be
a finitely generated subgroup of $G$. Then the inclusion of $H$
in $G$ is a quasi-isometric embedding if and only if $H$ is a
quasiconvex subgroup relative to the action of $G$ on $X$.
\e

Notes: Wise's, C.Hruska, has now proven this conjecture in the
case that $X$ is a CAT(0) $2$-complex. A counterexample, without the
isolated flats condition, is provided by $G=F_2\times\Z=<a,b>\times
<t>$ and $H=<at,bt>$ with the usual action of $G$ on
$(tree)\times\R$. $H$ is not quasi-convex since the intersection
$H\cap <a,b>$ is not finitely generated.

\section{Free Groups}

\b
(Hanna Neumann Conjecture) If $A$ and $B$ are nontrivial subgroups
of a free group, then $rk(A\cap B)-1\leq (rk(A)-1)(rk(B)-1)$. \e
Hanna
Neumann showed [Publ. Math. Debrecen 4 (1956), 186--189] $rk(A\cap
B)-1\leq 2(rk(A)-1)(rk(B)-1)$, and R.G. Burns [Math. Z. 119 (1971),
121--130] strengthened the inequality to $rk(A\cap B)-1\leq
2(rk(A)-1)(rk(B)-1)- \min(rk(A)-1,rk(B)-1)$. In many cases, the
conjecture holds, see Neumann, Walter D.: On intersections of finitely
generated subgroups of free groups. Groups---Canberra 1989, 161--170,
Lecture Notes in Math., 1456, Springer, Berlin, 1990. See also Dicks,
Warren: Equivalence of the strengthened Hanna Neumann conjecture and
the amalgamated graph conjecture. Invent. Math. 117 (1994), no. 3,
373--389.

\b
(Swarup) If $G$ is a Fuchsian group, define $area(G)$ to be the
area of the convex core of $\H^2/A$. Since $area(A)=2\pi
(rk(A)-1)$ for Fuchsian groups which are free, the Hanna Neumann
Conjecture can be phrased in terms of free Fuchsian groups: $2\pi
area(A\cap B)\leq area(A)area(B)$ whenever $A$ and $B$ are nontrivial
free subgroups
of a Fuchsian group. Prove such inequalities (possibly with a worse
constant) for (not necessarily free) torsion-free quasi-convex
subgroups of a quasi-convex Kleinian group in $\H^n$, where $area$
is replaced by the $n$-dimensional volume of the convex core. \e




\b
(J. Cornick) a) Let $G$ be a f.g. group such that every f.p. subgroup
is free. Is $G$ free?

b) If $G$ is f.g. and the homological dimension $hd\ G=1$, is $G$ free?
\e

Note: Yes to a) implies yes to b). There are counterexamples to a) in
which every f.p. subgroup is cyclic and torsion-free.



\subsection{Limit groups (Zlil Sela)}
The definition of limit groups involves group actions on $\R$-trees and is
given by Sela in his work on equations and inequalities over free groups. It
is a nontrivial theorem that a group $G$ is a limit group iff it is
$\omega$-residually free, i.e. for any finite subset $F$ of $G$ there is a
homomorphism from $G$ to a free group whose restriction to $F$ is
injective. Every limit group is finitely presented.

\b
{\bf Conjecture.} A finitely generated group admits a free action on an
$\R^n$-tree for some $n$ iff it is a limit group.\e

The ``if'' part follows from Sela's work.

\b
{\bf Conjecture.} Limit groups are $CAT(0)$.\e

\b
Characterize groups of the form $F_m*_\Z F_n$ which are limit
groups. \e

\b
Characterize 1-relator groups which are limit
groups. \e

\b (Swarup)
Show that a limit group is relative hyperbolic w.r.t its maximal 
abelian subgroups of rank bigger than one.\e

Sela showed that limit groups that contain no $\Z\times\Z$ are hyperbolic.

\subsection{Definable sets in $F_n$ (Zlil Sela)}

A subset $S$ of $F_n^k$ is {\it definable} if there is a first order sentence
with $k$ variables $x_1,x_2,\cdots,x_k$ which is true iff
$x_1,x_2,\cdots,x_k\in S$.

\b
{\bf Conjecture.} Let $n>1$. The set $$\{(x_1,x_2,\cdots,x_n)\in
F_n^n|x_1,x_2,\cdots,x_n \text{ is a basis of }F_n\}$$ is not definable.\e

\b
{\bf Conjecture.} The only definable subgroups of $F_n$ are cyclic and the
entire group. \e

\subsection{Genus in free groups (Zlil Sela)}
The {\it genus} of an element $x$ in the commutator subgroup of $F_n$ is the
smallest $g$ such that we may write $x=[y_1,y_2][y_3,y_4]\cdots
[y_{2g-1},y_{2g}]$ for some $y_i\in F_n$. Thus there is a map of the surface
of genus $g$ and one boundary component into the rose with boundary
corresponding to $x$. Applying an automorphism of the surface produces other
ways of writing $x$ as a product of $g$ commutators, and we call all such {\it
  Nielsen  equivalent}. One knows that there is a uniform bound $f(g)$
(independent of $x$) to the number of Nielsen equivalence classes in which an
element of genus $g$ can be written as a product of $g$ commutators.

\b
Give an explicit upper bound for $f(g)$.\e




\section{Transformation Groups}


\b
(de la Harpe) Is $PSL_2(\R)$ a maximal closed subgroup of
$Homeo_+(S^1)$? \e

\b Is there a proper closed subgroup of $Homeo_+(S^1)$ that acts transitively
on (unordered) 4-tuples? Or $k$-tuples ($k\neq 3$)? Relation to earthquakes?
\e

\b
(Christian Skau) There are many examples of groups acting on
non-homogeneous compacta with all orbits dense (boundaries of
word-hyperbolic groups, limit sets of Kleinian groups). Are there such
examples with group $=\Z$ (or amenable group)?\e

\b
(Swarup) Suppose $G$ is a 1-ended finitely presented group that
acts on a compact connected metric space $X$ as a convergence
group. What can be said about $G$ if $X$ has cut points? Does $X$ have
to be locally connected?  The model theorem of [Bowditch] and [Swarup]
says that if $G$ is word hyperbolic, then $X$ is locally connected and
doesn't have cut points.  In another interesting case, when $G$ is a
geometrically finite Kleinian group and $X$ its limit set, cut points
can arise, for example if $G$ splits over a parabolic subgroup.\e

\section{Kleinian Groups}


\b
(Jim Anderson) If $G$ is a group of isometries of $\H^n$, denote
by $Ax(G)$ the set of axes of the elements of $G$. If $G_1$ and $G_2$ are
finitely generated and discrete, does $Ax(G_1)=Ax(G_2)$ imply that $G_1$
and $G_2$ are commensurable?\e
 Yes if $n=2$ and $G_1,G_2$ are arithmetic by
[Long-Reid].

\b (G. Courtois, Gromov) Let $M=\H^n/\Gamma$ be a closed hyperbolic
manifold with $n\geq 3$. Does there exist a faithful non-discrete
representation $\rho:\Gamma\to Isom_+(\H^n)$ (same $n$)?  \e The
answer is yes (for certain manifolds) if the target is
$Isom_+(\H^{n+1})$ (Kapovich -- use bending deformations), and also by
Goldman's thesis if $n=2$. In the case of n=3 (Kapovich, Reid), if $M = {\bf
H}^3/\Gamma$ and the trace-field of $\Gamma$ admits a real galois
conjugate, then this can be used to induce a non-discrete faithful
representation of $\Gamma$ with real character. For example any
arithmetic Kleinian group derived from a quaternion algebra over a
field of degree at least $3$ has such a real representation.



\b
(Ed Taylor) Does there exist a constant $c>0$ such that the limit set
of every non-classical Schottky group has Hausdorff dimension $\geq c$.\e

\b
(Marden Conjecture) A torsion-free f.g. Kleinian group in dimension 3
is topologically tame, i.e. the quotient 3-manifold is the interior of
a compact 3-manifold.\e

\b
(Misha Kapovich)
Suppose that $G$ is a finitely generated Kleinian group in 
$Isom({\mathbb H}^n)$. Is it true that
$$
\delta(G) \le vcd(G)
$$      
with equality iff $G$ preserves a totally geodesic subspace 
${\mathbb H}^k\subset  {\mathbb H}^n$ so that  ${\mathbb H}^k/G$ is compact? 
\e

Here $\delta(G)$ is the exponent of convergence of $G$ and $vcd$ is the
virtual 
cohomological dimension. Note that the answer is positive for geometrically 
finite groups. 

A weaker form of this question is: 

\b (Misha Kapovich)
Is there $\epsilon >0$ so that 
$\delta(G)<\epsilon$ implies that $G$ is virtually free? \e

\b (Misha Kapovich)
Is there a finitely-generated discrete subgroup of $SO(n,1)$ whose 
action on the limit set is not ergodic? Is not recurrent? 
\e
Note that 
there are examples of finitely generated discrete subgroups of $SU(2,1)$ 
which do not act ergodically on the limit set (however in that example the 
action is recurrent). 


\b (Misha Kapovich -- a variation on a question of Bridson). 
Suppose that $G\subset 
SL(n,{\mathbb R})$ is a discrete subgroup of type $FP_{\infty}$ (over 
${\mathbb Z}$). Is it true that $G$ contains only finitely many $G$-conjugacy
classes of finite subgroups? \e

Note that there are examples (Feighn-Mess) where $G$ is a linear group 
of type $FP_k$ for arbitrarily large (but finite) $k$ so that $G$ contains 
infinitely many conjugacy classes of finite order elements. 

Update: This problem is completely solved (the answer is no) by Ian
Leary and Brita Nucinkis,

{\tt ftp://byrd.math.uga.edu/pub/html/archive/leary-nucinkis.html}

\section{Finiteness properties}

\b (Ian Leary)
Suppose $G$ is virtually of type FP over the field $F_p$ of $p$ 
elements, and let $g$ be an element of order $p$.  Is the 
centralizer of $g$ in $G$ also virtually of type FP over $F_p$?  
\e

If $G$ acts cocompactly on an $F_p$-acyclic space $X$, then 
the centralizer of $g$ acts cocompactly on $X^g$ = the set 
of fixed points of $g$, which is also $F_p$-acyclic by 
Smith theory.  So in this case the answer is yes.   

\b (Ian Leary)
Is there a group of finite vcd that does not act with finite 
stablizers on an acyclic complex of dimension equal to its 
vcd?  
\e

This question is sometimes called Brown's conjecture, 
although more often Brown's conjecture is taken to be the 
question of whether there is a universal proper G-space of 
dimension equal to the vcd.  The answer to this latter 
question is "no" --- Leary and Nucinkis have examples for each $n$ 
of groups for which the vcd is $2n$ but the minimal dimension
of the universal proper $G$-space is $3n$.

\b (Peter Kropholler)
If $G$ is FP over the rationals, is there a bound on the 
orders of finite subgroups of $G$?  
\e

Kropholler showed that this was the case for any $G$ that is 
both (a) of finite rational cohomological dimension and 
(b) type $FP_\infty$ over the integers.   

(Leary) This can be extended to the case when $G$ is only assumed to be 
$FP_n$ over the integers, for n = the rational coh. dim. of $G$.

\section{Splittings, Accessibility, JSJ Decompositions}


\b
(Potyagailo) {\rm Let $G\subset SO_+(n,1)$ be a geometrically finite,
torsion-free, non-elementary Kleinian group. Say that $G$ splits over
an essentially non-maximal parabolic subgroup $C\subset G$ as
$G=A*_{C}B$ (resp. $G=A*_C$) if $C$ is contained in a parabolic subgroup
$C'$ of higher rank such that $C'$ is not conjugate to a subgroup of
$A$ or $B$.}

Is the following true: $G$ is co-Hopfian iff it does not split over
trivial or essentially non-maximal parabolic subgroup? Note that the
latter condition is equivalent to saying that the parabolic splittings
are $K$-acylindrical (in the sense of Z.Sela) for a uniformly bounded
$K$.\e

\b
(Swarup) {\rm Let $G$ be a finitely presented group. Consider a maximal
graph of groups decomposition of $G$ with
finite edge groups and pass to the collection of vertex groups.
For each vertex group consider a maximal graph of groups decomposition
with 2-ended edge groups and pass to the collection of vertex groups.
Then split again along finite groups, then along two-ended groups  etc.}

{\bf Conjecture 1:} There is a finitely presented group for which this
process never terminates.

{\bf Conjecture 2 (Strong Accessibility):} For hyperbolic groups (and for
CAT(0) groups) this process always terminates.\e
Delzant-Potyagailo proved Strong Accessibility for hyperbolic groups without
2-torsion.

Note: For CAT(0) groups it would be natural to allow splittings over 
virtually abelian subgroup in the process. For general f.p. groups 
splittings over slender (small?) subgroups should be allowed.

\b
(Sageev) Is there a f.p. 1-ended group $G$ with $G\cong G*_{\Z}$? \e
Note
(Mitra) that such $G$ could not be co-Hopfian. In particular, $G$ could
not be torsion-free hyperbolic, by a theorem of Z. Sela.

\b
(Swarup) Show that the JSJ decomposition of a Haken 3-manifold $M$
depends only on $\pi_1M$ and not on the peripheral structure. Deduce
Johannson's deformation theorem: A homotopy equivalence $M\to N$
between Haken manifolds can be deformed so it induces a homeomorphism
between non-SFS components of the JSJ decompositions and induces a
homotopy equivalence between SFS pieces.\e


Update: This is answered by Scott-Swarup.

\b
(Papasoglou)
Is there a f.p. torsion-free group $G$ that does not split over a
virtually 
abelian subgroup, but has infinitely many splittings over $F_2$?
\e

This question is motivated by the fact that an irreducible atoroidal
closed 3-manifold has only finitely many incompressible surfaces of
any fixed genus.

\section{General Questions}




\b
(Eilenberg-Ganea) Is there a group $G$ of cohomological dimension 2 and
geometric dimension 3?\e

\b
(Shalen) If a finitely presented group $G$ acts nontrivially (i.e.
without global fixed points) on an $\R$-tree, does it act
nontrivially on a simplicial tree?\e

\b
(Mohan Ramachandran) {\rm See
[Napier-Ramachandran] for motivation. 
Consider the following two properties of a
finitely presented group $G$:

(A) $G$ virtually splits, i.e. some finite index subgroup of $G$ admits
a nontrivial action on a simplicial tree.

(B) Let $X$ be a finite complex with fundamantal group $G$. Then some
covering space of $X$ has at least two ends.

Most garden-variety groups satisfy both (A) and (B). Groups that
satisfy property (T) satisfy neither (A) nor (B).}

To what extent are (A)
and (B) equivalent? \e
The question makes sense for finitely generated groups
as well.

\b
(Bowditch) Let $\Gamma$ be the Cayley graph of an infinite
finitely generated group. Does there exist $K>0$ such that for all $R>0$
and all vertices $v\in\Gamma\setminus B(R)$ there is an infinite
ray from $v$ to $\infty$ which does not enter $B(R-K)$.\e

Anna Erschler observed that this fails for the lamplighter group.

\b
(Noel Brady) Are there groups of type $F_n$ but not $F_{n+1}$
($n\geq 3$) which do not contain $\Z\times \Z$? All known
examples contain $\Z^{n-1}$.\e





\b
(Olympia Talelli)
Is there a torsion-free group $G$ of infinite cohomological dimension such
that there is $n_0$ with the property that if $H$ is a subgroup of $G$ with
finite cohomological dimension $cd H$, then $cd H\leq n_0$.\e


\b
Can $\Z_{p^\infty}$ be embedded in an
$FP_\infty$-group? Or in an $FP_3$-group? Can an $F_n$-group be
embedded in an $F_{n+1}$-group ($n\geq 2$)? \e
Higman's celebrated embedding
theorem states that a group embeds into a finitely presented group iff it
has a recursive presentation. Any countable group embeds into a f.g. group. 


\b
Compute the asymptotic dimension of CAT(0) groups, $Out(F_n)$,
mapping class groups, nonuniform lattices, Thompson's group. Is there
a group of finite type whose asymptotic dimension is infinite? \e

Note (Misha Kapovich): Gromov states that the answer to the last question is
positive (and the group it does not admit a uniformly proper map into the
Hilbert space).

\subsection{Finite gap question (Jens Harlander)}
Let $G$ be a finitely generated group written as $G=F/N$ with $F$ a free group
of finite rank. Let $\Gamma$ be a finite graph with $\pi_1(\Gamma)=F$ and let
$\tilde\Gamma$ be the covering space of $\Gamma$ with $\pi_1(\tilde\Gamma)=N$.
Thus the deck group of $\tilde\Gamma$ is $G$. By $d_F(N)$ denote the smallest
number of $G$-orbits of 2-cells one needs to attach to $\tilde\Gamma$ to make
the space simply connected, and by $d_G(\frac{N}{[N,N]})$ denote the smallest
number of $G$-orbits of 2-cells one needs to attach to $\tilde\Gamma$ to kill
the first homology. Thus $d_F(N)<\infty$ iff $G$ is finitely presented and
$d_G(\frac{N}{[N,N]})<\infty$ iff $G$ is of type $FP_2$. There are examples
(Bestvina-N. Brady) of groups of type $FP_2$ which are not finitely
presented. A well-known question, stemming from the 1965 work off C.T.C. Wall,
becomes:

\b Is there an example of a {\it finitely presented} group $G=F/N$ such that
   $d_G(\frac{N}{[N,N]})<d_F(N)(<\infty)$? \e 
   
   One approach to construct such an example is the following. Suppose $H=F/N$
   is finitely presented and contains as subgroups groups of the form
   $C^n=C\times C\times \cdots\times C$ for all $n$. Define $G_n=H*_{C^n}H$
   which is written as $F*F/N_n$ in the obvious way. If $C$ is finite and
   perfect, then $d_{F*F}(\frac{N_n}{[N_n,N_n]})$ is independent of $n$. 

\b If $C$ is nontrivial and finite, is $\lim_{n\to\infty}d_{F*F}(N_n)=\infty$?
\e 


\section{2-complexes}

\b (Whitehead) Is every subcomplex of an aspherical 2-complex aspherical?\e

\b (Andrews-Curtis) If $K$ and $L$ are simple homotopy equivalent
finite 2-complexes, can one transform $K$ to $L$ by a sequence of elementary
collapses and expansions of 1- and 2-cells, and by sliding 2-cells
(i.e. reattaching them by maps homotopic to the old attaching maps)?\e
This is unknown even when $K$ is contractible and $L$ is a point.

\b
(Wise)
Is there a finite aspherical 2-complex $X$ with $\pi_1(X)$ coherent
and with $\chi(X)\geq 2$?
\e

\section{Symmetric Spaces}


\b (Igor Belegredek)
Let $X$ be a non-positively curved symmetric space.
Find conditions on a group $\Gamma$ so that the space of conjugacy
classes of faithful discrete representations of $\Gamma$
into the isometry group of $X$ is compact (non-compact).\e

Notes (Belegradek): 
Suppose a sequence of faithful discrete representations goes to infinity. 
That defines a natural isometric action of $\Gamma$ on the asymptotic
cone of $X$, that is an Euclidean Tits building according to Kleiner-Leeb.
The action is small (in a certain sense) and has no global fixed point. 
Thus it is enough to understand what groups cannot (can) have such action.
First example to look at is when $X$ is a product of two rank one spaces,
so the asymptotic cone of $X$ is a product of two trees.


\b
(Belegredek)
Is there a 3-complex $X$ (not necessarily aspherical) which is
not homotopy equivalent to a 2-complex but $H^3(X;\{G\})=0$ for all
local coefficients? \e

The condition is equivalent to saying that $X$ is dominated by a
finite 2-complex, or that $id:X\to X$ is homotopic to a map into the
2-skeleton. 

\section{Quasi-Isometries}

\b
Study the quasi-isometry group $QI(\R^n)$. How big is it?\e

\b
(Kleiner) What are the quasi-isometries of the 3-dimensional
group $Sol$?\e

\b
(Kleiner) What are the q.i's of the Gromov-Thurston examples
of negatively pinched manifolds?\e

\b
(Feighn) {\rm Let $\phi:F_n\to F_n$ be an automorphism of the free
group $F_n$ and let $M_\phi$ be its mapping torus. }

Classify these groups up to quasi-isometry.\e

Note: It was shown by [Macura] that if two automorphisms have polynomial
growths with distinct degrees, then their mapping tori are not
quasi-isometric. 

\b
(Bridson) Same as above for automorphisms $\Z^n\to \Z^n$.\e

\b
(Bridson) Is $G\times \Z$ 
quasi-isometric to $G$ for $G=$Thompson's group? Any
f.g. group?\e






\section{Mapping Class Groups, $Out(F_n)$}

\b
(Levitt) Can every measured geodesic lamination with 2-sided
leaves on a non-orientable compact hyperbolic surface be approximated
by a simplicial measured geodesic lamination with 2-sided leaves?\e
Remark: The subspace $\Omega$ of the sphere of geodesic measured
laminations that contain 1-sided leaves is open and dense (of full
measure, [Danthony-Nogueira]). More generally, understand the action
of the mapping class group on the space of projectivized measured
laminations. Is the action minimal on the complement of $\Omega$.  

\b
(Lubotzky) Does $Out(F_n)$ have the congruence subgroup property?\e
If $G\subset F_n$ is a characteristic subgroup of finite index, then kernels
of homomorphisms to $Out(F_n/G)$ are called congruence subgroups.

\b
(Kapovich) a) Let $\Sigma_g$ be the closed orientable surface of
genus $g$. Is there a faithful representation $\pi_1(\Sigma_g)\to
MCG(\Sigma_h)$ into a mapping class group such that the image consists
of pseudo-Anosov classes plus identity (for some $g,h>1$)?

b) Is there a 4-manifold $M$ which is a surface bundle over a surface
s.t.  $\pi_1(M)$ is word-hyperbolic? Or so that $M$ is hyperbolic?\e
                                                                       
\b
(Bowditch) Is the Weil-Petersson metric on Teichm\" uller space
hyperbolic?  Is it quasi-isometric to the curve complex? \e
The latter is
word-hyperbolic by [Mazur-Minsky]. 

Update: J. Brock proved that W-P space is quasi-isometric to the
``pants complex'' ({\tt
http://www.math.uchicago.edu/$\sim$brock/home/text/papers/wp/www/wp.ps.gz}).

Brock-Farb have announced (October 2000) that the space is
hyperbolic for a torus with $\leq 2$ punctures and a sphere with $\leq
5$ punctures, but not hyperbolic for more complicated surfaces. A
related result was announced by Sumio Yamada: the completion of the
Weil-Petersson metric is a CAT(0) space with ideal points coming in
natural pieces, with each piece the metric product of simpler
Weil-Petersson spaces.

\b
(Brock) What is the rank of the Weil-Petersson metric? What is the
rank of the mapping class group?
\e

The {\it rank} of a metric space $X$ is the maximal $n$ such that $X$
admits a quasi-isometric embedding $\R^n\to X$. Brock-Farb show that a
lower bound for W-P is $\frac 12 (3g-3+b)$ (Yamada's result above predicts
that the rank is the maximal number of disjoint subsurfaces with
noncompact Teichm\" uller space, and the Brock-Farb result confirms
this as a lower bound). For the mapping class group, a lower bound is
given by a maximal curve system -- the associated group of Dehn twists
is embedded quasi-isometrically.
                                                                      

\subsection{Automorphisms of free groups (Gilbert Levitt)}
$\alpha:F_n\to F_n$ will denote an automorphism. 

\b
Assuming that the fixed subgroup $Fix(\alpha)$ is cyclic, find a bound on
the length of a generator of $Fix(\alpha)$ in terms of the complexity of
$\alpha$. \e
(ed. comm.: an easier version of the question would be to bound the
length of the generator of $Fix(\alpha)$ in terms of the complexity of a
relative train-track representative of $\alpha$.)

\b
Which $\alpha$ preserve an order (invariant under right translations) on
$F_n$? If $\alpha$ has periodic elements it cannot preserve an order. Are
there other obstructions? \e
A result of Perron-Rolfsen asserts that if
$\alpha_{ab}\in GL(n,\Z)$ has all eigenvalues $>0$ then $\alpha$ preserves an
order. 

\b
Does $Out(F_n)$ ($n>2$) have a right orderable subgroup of finite index?\e
It is a theorem of Dave Witte (Proc.A.M.S. 122 (1994), no. 2, 333--340) that
$SL_n(\Z)$ ($n>2$) does not have a right orderable subgroup of finite
index. Mapping class groups of surfaces with nonempty boundary are orderable
([Short-Wiest] following Thurston). 

\b Can the mapping class group (or a finite index subgroup of it) of a closed
surface be embedded into $Out(F_n)$? Into the mapping class group of a
punctured surface?  \e


\b Do finite index subgroups of $Out(F_n)$ ($n>2$) have property $(FA)$, i.e.
does every isometric action of such a subgroup on a simplicial tree (or even
an $\R$-tree) have a global fixed point?\e

Update: Lubotzky constructed a finite index subgroup of $Out(F_3)$
that maps onto a nonabelian free group.


\b
(Grigorchuk)
$Aut(F_n)$ acts on the space $\partial_3 F_n$ of triples of distinct
ends of $F_n$. Denote by $Y_n$ the compact space (Cantor set) the
quotient space of $\partial_3F_n$ by the group of inner
automorphisms. Thus $Out(F_n)$ acts on $Y_n$. Describe the dynamics of
this action; in particular the dynamics of any individual outer
automorphism. 
\e

\subsection{Schottky groups in mapping class groups (Lee Mosher)}
Recall that a {\it Schottky group} is a subgroup $F$ of $Isom(\H^n)$ which is
free of finite rank, discrete, consists of loxodromic elements, and  every
orbit is quasi-convex.  Similarly, we define a Schottky subgroup of
$Isom(Teich(S))\cong MCG(S)$ where we replace ``loxodromic'' by
``pseudoAnosov''. Note that $MCG(S)$ acts discretely on $Teich(S)$, so the
word ``discrete'' can be omitted from the definition.

\b
Do there exist f.g. free subgroups of $MCG(S)$ consisting of identity and
   pseudoAnosov mapping classes which are not Schottky?\e

If such a subgroup $F$ exists, then $\pi_1(S)\rtimes F$ is not
word-hyperbolic, but has no Baumslag-Solitar subgroups. Potential candidates
for such $F$ might be found as subgroups of Whittlesey's example.

\b Do there exist non-free pseudoAnosov subgroups?\e

\b Is there an $Out(F_n)$ analog of the above? \e

One difficulty is that the geometry of Outer Space is not well understood. It
is still true (Bridson-Vogtmann) that the group of simplicial homeomorphisms
of Outer Space is $Out(F_n)$. It is not clear what a good metric on Outer
Space should be. For example, if $\alpha$ is an irreducible automorphism, then
the set of train-tracks for $\alpha$ should be a quasi-line, and these
quasi-lines for $\alpha$ and $\alpha^{-1}$ should be a uniform distance apart
(in the case of $MCG$ they coincide). Our ignorance is exemplified in the
following:

\b Given $n$ is there a uniform constant $K$ such that for every irreducible
   automorphism $\alpha:F_n\to F_n$
   $$\frac{\log\lambda(\alpha)}{\log\lambda(\alpha^{-1})}\leq K$$ where
   $\lambda(f)$ denotes the growth rate of $f$?\e ($K=1$ in the case of $MCG$.)

\subsection{Betti numbers of finite covers (Andrew Casson)}

\b \label{1}
Does every automorphism $h:F_n\to  F_n$ leave invariant a finite index
   subgroup $K$ such that $h_{ab}:K/K'\to K/K'$ has an eigenvalue which is a
   root of unity?\e

\b \label{2}
Does every closed 3-manifold $M$ which fibers over $S^1$ with fiber of
   genus $\geq 2$ have a finite cover $\tilde M$ with $b_1(\tilde M)>1$?\e
   
   Note: Q \ref{2} is a ``warmup'' for the well-known conjecture that every
   aspherical 3-manifold has a finite cover with positive $b_1$. Q \ref{2} is
   Q \ref{1} with surface groups instead of free groups.

\subsection{QI rigidity of MCG and $Out(F_n)$ (Martin Bridson)}

\b Is $MCG(S_g)\to QI(MCG(S_g))$ an isomorphism for $g\geq 3$? \e
For $g=2$ the
hyper-elliptic involution is central and has to be quotiented out from the
left-hand side.

\b Suppose that $\phi:MCG(S_g)\to MCG(S_g)$ is a quasi-isometry. Does $\phi$
   map maximal flats to maximal flats (``maximal flats'' come from maximal
   rank abelian subgroups)?\e

Note: By work of Ivanov, yes to \#2 implies yes to \#1.

\b Same questions for $Out(F_n)$.\e

\b If $\Gamma$ is an irreducible uniform lattice in a higher rank connected
   semisimple Lie group, does every homomorphism $\Gamma\to Out(F_n)$
   necessarily have finite image?\e

This is true for nonuniform lattices as observed by Bridson-Farb (consequence
of the Kazhdan-Margulis superrigidity and the fact that solvable subgroups of
$Out(F_n)$ are virtually abelian [Bestvina-Feighn-Handel]).  


\subsection{Linearity of mapping class groups (Joan Birman)}
\b Is $MCG(S_{g,b,n})$ linear? \e
Here $g$ is the genus, $b$ the number of
boundary components and $n$ the number of punctures. By the work of Bigelow
and Krammer, $MCG(0,1,n)$, $MCG(0,0,n)$, and $MCG(2,0,0)$ are linear.


\subsection{The Singular Braid Monoid (Joan Birman)}
A {\it singular braid} is defined by a braid diagram as usual, except that in
addition to over- and under-crossings, crossings of strands are allowed. The
standard singular braids are
\begin{itemize}
\item braid $\sigma_i$ in which strand $i+1$ goes over the $i^{th}$ strand,
  and
\item singular braid $\tau_i$ in which strands $i$ and $i+1$ cross.
\end{itemize}

The set of singular braids on $n$ strands forms a monoid $SB_n$ generated by
$\{\sigma_1^{\pm 1},\cdots,\sigma_{n-1}^{\pm 1},\tau_1,\cdots,\tau_{n-1}\}$. A monoid
presentation of $SB_n$ is given in terms of these generators and the following
relations:
\begin{itemize}
\item braid relations $\sigma_i\sigma_j=\sigma_j\sigma_i$ for $|i-j|>1$ and
  $\sigma_i\sigma_{i+1}\sigma_i=\sigma_{i+1}\sigma_i\sigma_{i+1}$,
\item $\tau_i\tau_{j}=\tau_{j}\tau_i$ for $|i-j|>1$,
\item $\sigma_i\tau_j = \tau_j\sigma_i$ if $i=j$ or $|i-j|>1$,
\item $\tau_i\sigma_{i+1}\sigma_i=\sigma_{i+1}\sigma_i\tau_{i+1}$,
\item $\tau_i\sigma_{i+1}\sigma_i=\sigma_{i+1}\sigma_i\tau_{i+1}$.
\end{itemize}

See J. Birman: "New points of View in knot
theory", Bull AMS {\bf 28}(1993), no. 2, 253-287.

The singular braid monoid is closely related to Vasiliev invariants of knots.
Define a monoid homomorphism $\Phi:SB_n\to\Z B_n$ to the group ring of the
braid group $B_n$ by $\Phi(\sigma_i)=\sigma_i$ and
$\Phi(\tau_i)=\sigma_i-\sigma_i^{-1}$.

\b
{\bf Conjecture.} $\Phi$ is injective.\e

The proof that this is true for $n=3$ is
due to A.Jarai, "On the monoid of singular braids", Topology and its
Applications, {\bf 96}(1999), 109-119.

It is also true on singular braids with at most two
singularities (exercise). 


\section{Other Questions}

\b (Kevin Whyte)
Let $K$ be a finite complex with $\pi=\pi(K)$ amenable. Is there a uniform
bound to the betti numbers of finite covers of $K$? \e
Related to the work of
Andrzej Zuk on $\ell_2$-cohomology. Also related to the conjecture that such
$\pi$ are elementary amenable.

\b (Kevin Whyte)
Is every solvable $PD(n)$ group polycyclic?\e


\b
(Henry Glover) Does for every finite graph $G$ the following
$1-2-\infty$-conjecture hold: $G$ is planar, a double cover of $G$ is
planar, or no finite cover is planar? \e
There is a reduction (ref??) to
the graph $G$ obtained by coning off the 1-skeleton of the octahedron.

\b
(S. Ivanov) Is there a f.p. slender group which is not
polycyclic-by-finite? \e
A group is slender (or Noetherian) if every
subgroup is f.g. [Olshanskii] has constructed a f.g. counterexample.




\b (Seymour Bachmut)
Is $SL_2(K)$ finitely generated for $K=\Z[X,X^{-1}]$ or
$K=F[X,X^{-1},Y,Y^{-1}]$ for a field $F$?\e

\b Are 1-relator groups coherent?\e

\begin{comment}
\b
(Y. Shalom)
Is the following true: If $\Gamma_1,\Gamma_2$ are two lattices in a
semisimple Lie group $G$ and $\Q-rank(\Gamma_1)<\Q-rank(\Gamma_2)$
then $\Gamma_1$ does not embed uniformly into $\Gamma_2$.
\e
\end{comment}

\subsection{Hopf-Thurston Conjecture (Mike Davis)}

The Hopf-Thurston Conjecture asserts that if $M$ is a closed aspherical
$2n$-manifold then $(-1)^n\chi(M)\geq 0$. The stronger Singer Conjecture
asserts that the $\ell_2$-homology of $\tilde M$ is 0 except possibly in
dimension $n$. See the recent Davis-Okun work. A special case (using
right-angled Coxeter groups) of Hopf-Thurston's Conjecture is the following:

Let $L$ be a flag triangulation of $S^{2k-1}$ and define
$$\chi=1-\frac{1}{2}f_0+\frac{1}{4}f_1-\frac{1}{8}f_2+\cdots=1-\sum_{i=0}^{2k-1}(-1)^i\frac{1}{2^{i+1}}f_i$$
  where $f_i$ is the number of $i$-simplices of $L$.

\b
{\bf Conjecture.} $(-1)^k\chi\geq 0$.\e

This is true for $S^3$ (Davis-Okun) and it is also true if $L$ is the
barycentric subdivision of another triangulation. 




                                                            
\subsection{Word Problem (Martin Bridson)}
\b Do there exist groups $G$ with balanced presentation (same number of
   generators and relations), with $H_1(G)=0$ and with unsolvable word
   problem?\e

Note: The standard examples of groups with unsolvable word problem have more
relations than generators. The condition $H_1(G)=0$ is added to rule out
counterexamples obtained by adding silly generators. Any other condition that
rules this out is acceptable.

\b Is there a sequence of (perfect, of course) groups with balanced
presentations among which one cannot recognize trivial groups?  \e   


\subsection{Membership Problem in semigroups (John Meakin)}
Let $G$ be a 1-relator group with the relator $W$ a cyclically reduced word in
the generators. Let $P$ be the submonoid of $G$ generated by all prefixes of
$W$. 

\b Is the membership problem for $P$ in $G$ decidable?\e

For motivation and special cases, see [S. Ivanov - S. Margolis - J. Meakin,
One relator inverse monoids and 1-relator groups, J. Pure Applied Alg.]. The
paper reduces the word problem in $G$ to the question above.


\end{document}
